\documentclass[11pt,a4paper]{article}
\usepackage[margin=1in]{geometry}
\usepackage{hyperref}
\usepackage{enumitem}
\usepackage{graphicx}
\usepackage{xcolor}

\makeatletter
\newcommand{\BvrSetLabInstructor}[1]{\gdef\Bvr@LabInstructor{#1}}
\newcommand{\Bvr@LabInstructor}{Dr. Raghav B. Venkataramaiyer}
\newcommand{\BvrLabInstructorValue}{\Bvr@LabInstructor}
\makeatother
\usepackage{titling}
\pretitle{\begin{center}\bfseries\fontsize{18bp}{18bp}\selectfont}
\makeatletter
\newcommand{\BvrSubtitle}[1]{\posttitle{\par\end{center}\begin{center}\large#1\end{center}\vskip 0.5em}}
\makeatother
\newcommand{\BvrHint}[1]{{\normalfont\normalsize\itshape\hfill #1}}

\title{EcoCode: Green-Software Analysis Platform}
\BvrSetLabInstructor{Dr. Raghav B. Venkataramaiyer}
\BvrSubtitle{%
  Project Proposal for UCS503P \\
  Submitted to: \BvrLabInstructorValue \\ \bfseries
  Thapar Institute of Engineering and Technology
}
\author{%
  Gurnoor Singh Sodhi\thanks{Roll No: \texttt{1024030387}}%
  \and Rakesh\thanks{Roll No: \texttt{1024030393}}%
  \and Manav\thanks{Roll No: \texttt{1024030398}}%
}
\date{September 2026}

\begin{document}
\maketitle

\section{Project Idea}
EcoCode is a green-software analysis platform for Java projects. It is intended to analyse complete multi-file repositories, identify source-code patterns that may lead to unnecessary computation or resource usage, and provide evidence-based recommendations. The analysis is deterministic and rule-based rather than AI-dependent.

\section{Problem Statement}
Software can waste CPU time, memory, I/O, database operations and developer effort because inefficient patterns are difficult to identify consistently in large code bases. Existing manual review is time-consuming and may miss cross-file relationships. EcoCode aims to automate the first stage of this review by building a project model, extracting method relationships and detecting potential resource-wasting patterns.

\section{Proposed Solution}
The system will incrementally provide:
\begin{enumerate}[leftmargin=*]
  \item whole-project Java indexing and AST parsing;
  \item class, method and method-call modelling;
  \item lightweight symbol resolution and call-graph construction;
  \item deterministic static-analysis rules (EC101--EC108);
  \item later-stage build execution, runtime profiling, evidence validation, recommendations and version comparison.
\end{enumerate}

\section{Development Methodology}
The project follows an Agile iterative and incremental model. Early iterations deliver a usable static-analysis pipeline before runtime intelligence and product integration are added. This allows each stage to be tested on real open-source repositories and refined before the next stage.

\section{Initial Technology Stack}
\begin{itemize}[leftmargin=*]
  \item Java 21 and Gradle
  \item JavaParser for AST analysis
  \item JUnit 5 for automated tests
  \item Planned later-stage tools: Docker, JFR/JMC, JGit and a web dashboard
\end{itemize}

\section{Expected Outcome}
A working prototype that can accept a Java project path, index and parse the repository, construct a project model and call graph, execute deterministic rules, and report source-level static findings. Rule accuracy will be validated and refined using multiple repositories before being treated as final recommendations.

\end{document}
